\documentclass[11pt]{article}
\usepackage[T1]{fontenc}
\usepackage{lmodern,microtype,amsmath,amssymb,amsthm,mathtools,booktabs}
\usepackage[margin=1.05in,headheight=14pt]{geometry}
\usepackage{xcolor,fancyhdr}
\usepackage[colorlinks=true,linkcolor=teal!65!black,citecolor=teal!65!black,urlcolor=teal!65!black,pdfauthor={OpenAI ChatGPT},pdftitle={Recovering Dirichlet characters from short windows of zeros}]{hyperref}
\newtheorem{theorem}{Theorem}[section]
\newtheorem{lemma}[theorem]{Lemma}
\newtheorem{corollary}[theorem]{Corollary}
\newtheorem{proposition}[theorem]{Proposition}
\newcommand{\R}{\mathbb R}
\newcommand{\C}{\mathbb C}
\DeclareMathOperator{\sinc}{sinc}
\newcommand{\FQ}{\mathcal F(Q)}
\newcommand{\Lthree}{\mathcal L_3}
\pagestyle{fancy}\fancyhf{}\fancyhead[L]{\small Dirichlet characters from zero ordinates}\fancyhead[R]{\small Research note}\fancyfoot[C]{\thepage}
\setlength{\parskip}{2pt}\emergencystretch=1.3em
\begin{document}
\begin{center}
{\small\sffamily\color{teal!70!black} ANALYTIC NUMBER THEORY}\par\vspace{9pt}
{\LARGE\bfseries Recovering Dirichlet characters\\[3pt]from short windows of zeros}\par\vspace{9pt}
{\normalsize November 27, 2026}
\end{center}
\begin{abstract}
For every $Q\ge3$, the positive ordinates up to height
$O(\log^2(2Q)/\log\log(3Q))$, counted with multiplicity, distinguish all primitive
Dirichlet characters of conductor at most $Q$. The result includes conjugate complex
characters and the conductor-one character. It uses Dirichlet GRH in its critical-line
form. A translated window centered at $U$ works with half-length enlarged by
$b\log\log\log |U|$ for any $b>1/\pi$ and large $|U|$. The proof uses a smooth
explicit formula together with an upper-bound sieve for pairs of primes.
Many small primes separate each pair of characters, so individual nearby prime powers
can be tolerated on average. We give a quantitative lower bound for the cost of
matching the two lists of ordinates and a finite-alphabet version of the recovery.
\end{abstract}
\section{The identification theorem}
Write $\FQ$ for the primitive Dirichlet characters of conductor at most $Q$, including
the primitive principal character of conductor one. All logarithms are natural. Set
\[
 W(Q)=\frac{\log^2(2Q)}{\log\log(3Q)},\qquad
 \Lthree(t)=\log\log\log(e^{e^e}+t)\quad(t\ge0).
\]
For an interval $I$ let $\Gamma_\chi(I)$ be the multiset of real $\gamma\in I$ for
which $L(1/2+i\gamma,\chi)=0$, with orders as multiplicities. Dirichlet GRH identifies
these with the imaginary parts of all nontrivial zeros in the interval.
\begin{theorem}[Uniform identification]\label{thm:main}
Use Dirichlet GRH. For each $b>1/\pi$ there is $C_b$ such that, for all $Q\ge3$
and $U\in\R$, the map $\chi\mapsto\Gamma_\chi([U-H,U+H])$ is injective on $\FQ$ if
\begin{equation}\label{eq:H}
 H\ge C_b W(Q)+b\Lthree(|U|).
\end{equation}
In particular there is an absolute constant $C$ such that
$\chi\mapsto\Gamma_\chi((0,CW(Q)])$ is injective.
\end{theorem}
The last assertion distinguishes $\chi$ from $\bar\chi$ using a positive-height
interval. Folding a symmetric list by $\gamma\mapsto|\gamma|$ discards this
information because conjugation reflects ordinates. Each window in the theorem needs
only its own zeros. Neither conductor nor parity is an input to the measurements.
Primitivity is essential: inducing a character to a multiple of its modulus changes
Euler factors whose zeros lie on $\Re s=0$ and preserves its open-strip zeros.

\paragraph{Relation to the two repositories.}
The formal input is \texttt{LFunction\_eq\_zero\_re} in the readable edition
of \cite{lean}. Its statement in Mathlib's language covers all nonzero moduli.
We inspected the exported statements and their definitions in the pinned source.
The upstream authors report kernel and independent checker runs with the usual
three Lean axioms. This paper uses their theorem and supplies a separate analytic
proof of Theorem~\ref{thm:main}. We did not rebuild the upstream project or formalize
the present proof.

The merged ProveIt article \cite{article} describes itself as a synthesis of 103 reports.
Its result \texttt{rec:cor:quadsep} recovers primitive nonprincipal quadratic
characters from both signs up to $O(\log^2(2Q)\log\log(3Q))$. Its question
\texttt{rec:sec:open}, item~2, asks to remove the extra logarithm and to treat complex
characters. The two cutoff bounds differ by a factor of order $(\log\log Q)^2$.
The article also gives a triple-logarithmic window for fixed families
(\texttt{zer:thm:short} and \texttt{rec:thm:localrank}). Our contribution in the
translated form is the uniform cost $C_b W(Q)$ for the whole family of conductors.
Its fixed-family signed-excess result gives information about each direction of
multiplicity difference that injectivity alone does not recover.

\paragraph{Comparison with the literature.}
Fujii obtained positive proportions of multiplicity discrepancies for fixed characters
\cite{fujii}. Banks proved an unconditional uniform short-interval discrepancy theorem
for equal moduli with width $O_\theta(q^\theta\log T)$ when $T\gg_\theta q^\theta$,
$\theta>1/3$, improved to $\theta>1/4$ for cubefree $q$ \cite{banks}. Banks and
Loftus prove existence of genuinely unshared zeros for fixed characters in sufficiently
high windows in a different size range \cite{loftus}. Our statistic detects
differences of orders, including at a common zero. This comparison establishes a
specific advance over the linked article. A targeted search of related work did
not locate the uniform low-height bound of Theorem~\ref{thm:main}. Publication
priority remains to be assessed independently.

\section{Many primes distinguish every pair}
The arithmetic input is stronger than the existence of one small prime at which
the characters differ. A positive fraction of the primes in one common interval
give an order-one separation, uniformly in the pair.
\begin{lemma}\label{lem:witness}
There are absolute $A_0,c_0>0$ such that, for $A\ge A_0$,
$x=A\log^2(2Q)$ and distinct $\chi_1,\chi_2\in\FQ$, at least $c_0 x/\log x$
primes $x<p<2x$ are prime to both conductors and satisfy
$|\chi_1(p)-\chi_2(p)|\ge1$.
\end{lemma}
\begin{proof}
Choose once and for all a smooth nonnegative $w$ supported in $(1,2)$ with
$\int w(y)\,dy=1$. Let $M$ be the least common multiple of the conductors.
The quotient character $\chi_1\overline{\chi_2}$ on the units modulo $M$ has a
nonprincipal primitive inducer $\psi$ of conductor $q\le M\le Q^2$. Indeed, equality
of the two induced characters on the units would give them the same primitive inducer.

For this $\psi$ the smooth explicit formula and Dirichlet GRH give
\begin{equation}\label{eq:smoothed}
 \sum_{n\ge2}\Lambda(n)\psi(n)w(n/x)\ll_w \sqrt{x}\log(2q).
\end{equation}
For detail, use $g_x(u)=e^{u/2}w(e^u/x)$ in the explicit formula of
Section~\ref{sec:localized-kernel}. Its positive prime terms give exactly the sum.
Its transform is $\sqrt{x}x^{it}\int w(y)y^{-1/2+it}\,dy$, of size
$O_w(\sqrt{x}(1+|t|)^{-3})$ by integration by parts. The number of ordinates
in any unit interval near $t$ is $O(\log(2q(2+|t|)))$ by the classical zero
count \cite{bmor}, also for either sign by taking an upper bound from the two-sided
count. The archimedean integral is $O_w(1)$ because $u\ge\log x$ on the support
and the density is $O(e^{-u/2})$. Both the conductor term and the pole term vanish.

The total weighted prime mass in this interval is $(1+o_w(1))x$ by the prime
number theorem. In \eqref{eq:smoothed} removing higher powers costs
$O_w(\sqrt{x}\log^2(2x))$ by elementary counting and removing the primes of
$M$ costs $O_w(\log(2M))$. Choose $A_0$ large enough that the remaining
prime mass is $\ge 3x/4$ and its character-weighted sum has modulus $\le x/8$.
The implied constants are uniform since $\log(2M)\ll\log(2Q)$.
If $S$ denotes the weighted mass of primes with $\Re\psi(p)\le1/2$, the real
part of that sum is at least $\tfrac12(3x/4)-\tfrac32 S$.
Thus $S\ge x/6$. Each prime carries $O_w(\log x)$ weight and at every prime
so counted $|\chi_1(p)-\chi_2(p)|^2=2-2\Re\psi(p)\ge1$.
\end{proof}

\section{A localized explicit formula}
\label{sec:localized-kernel}

Fix $b>1/\pi$ and $1/2<d<\pi b/2$ (one may fix the midpoint).
Use $\sinc v=\sin(v)/v$ with value $1$ at zero. Put
\[
 k(t)=e^{-\cosh(t/b)},\qquad
 \eta(v)=\frac{\int_{\R}k(t)e^{-itv}\,dt}{\int_{\R}k(t)\,dt},
 \qquad K_B(v)=\eta(v)\sinc^4(Bv).
\]
For a prime $p$, write $\ell=\log p$. For $B\ge1$ and $U\in\R$ define
\[
 g(u)=\frac{u}{\ell}e^{-iU(u-\ell)}K_B(u-\ell),\quad
 h(t)=\int_{\R}g(u)e^{itu}\,du,\quad
 D_{p,B,U,H}(\chi)=-\frac{\sqrt p}{\ell}
       \sum_{|\gamma-U|\le H}h(\gamma).
\]
The ordinates are signed and counted with multiplicity. All constants below
may depend on $b,d$. Write $\delta_\chi=1$ for the conductor-one character and
$0$ otherwise. Define the nonnegative leakage from nearby prime powers by
\[
 L_B(p)=\sum_{\substack{p/e\le n\le ep,\ n\ne p\\ \Lambda(n)>0}}
   \left(1+B\frac{|n-p|}{p}\right)^{-4}.
\]

\begin{lemma}[Prime-value decoder]\label{lem:kernel-decoder}
For primitive $\chi$ of conductor $q\le Q$, under Dirichlet GRH, and for
$H=4B+\Delta$, $\Delta\ge1$, one has
\begin{align*}
 |D_{p,B,U,H}(\chi)-\chi(p)|\ll_{b,d}{}&L_B(p)+\frac{p}{B^4}+\frac1B
       +\delta_\chi\frac{p}{B\ell}\\[-2pt]
 &+\frac{\sqrt p}{\ell}\log\!\big(eQ(3+|U|+H)\big)
                        \exp\{-c_b e^{\Delta/b}\}
\end{align*}
with some $c_b>0$. Given $\epsilon>0$, the last term, with its implied
constant, is $\le\epsilon$ whenever
\begin{equation}\label{eq:kernel-tail-choice}
 \Delta\ge C_{b,d,\epsilon}+b\log\!\big(2+\log(2p)+
           \log\log(e^e Q(3+|U|+B+p))\big).
\end{equation}
The normalized summand $-\sqrt p\,h(t)/\ell$ is bounded by
$O_{b,d}(\sqrt p/(B\ell))$ and has Lipschitz constant
$O_{b,d}(\sqrt p/B)$ on $\R$.
\end{lemma}

\begin{proof}
The defining integral gives $|\eta(v)|\le1$. Shifting its contour by
$-id\operatorname{sgn} v$ yields $|\eta(v)|\ll_{b,d}e^{-d|v|}$, since
$|k(t\pm id)|=\exp(-\cos(d/b)\cosh(t/b))$. Thus $h$ is analytic beyond
the closed strip $|\operatorname{Im}t|\le1/2$. In our Fourier convention,
the Guinand--Weil formula \cite[Eqs.~(2.1) and (2.6)]{weil} reads, when $g(0)=0$,
\begin{align*}
 \sum_\gamma h(\gamma)={}&\delta_\chi\{h(i/2)+h(-i/2)\}\\
 &-\int_0^\infty \frac{e^{-(a+1/2)u}}{1-e^{-2u}}
                 \{g(u)+g(-u)\}\,du\\
 &-\sum_{n\ge2}\frac{\Lambda(n)}{\sqrt n}
            \{\chi(n)g(\log n)+\overline{\chi(n)}g(-\log n)\},
\end{align*}
where $a\in\{0,1\}$ is the parity. For the sign convention, Mellin inversion
of $\int g(u)e^{(s-1/2)u}\,du$ against $-L'/L(s,\chi)$ on the right of $1$
selects $g(\log n)$. On the reflected line the functional equation supplies
$\overline\chi$ and $g(-\log n)$. The conductor term is a multiple of
$g(0)$, and the integral follows from the digamma integral by substituting
$t=2u$. Admissibility follows since $k$ and all its derivatives decrease
doubly exponentially in each smaller horizontal strip
$|\operatorname{Im}t|<\pi b/2$, and the transform of $K_B$ is an average of
translates of $k$ over a compact interval, as shown below.

We have $g(\ell)=1$. In the positive-frequency sum with
$|\log(n/p)|\le1$, the quantities $\sqrt{p/n}$, $|\log n|/\ell$ and
$\Lambda(n)/\ell$ are bounded absolutely, and
$|\log(n/p)|\ge |n-p|/(ep)$. Hence the normalized error there is
$O(L_B(p))$. For
$\log(n/p)\in[j,j+1)$, $j\ge1$, there are $O(pe^j)$ integers and each
normalized term is
$O_{b,d}((1+j)^2 e^{-(d+1/2)j}/B^4)$. Their sum is $O_{b,d}(p/B^4)$.
For $n<p/e$, the ratios $(\log n)/\ell$ and $\Lambda(n)/\ell$ are at most
one and $\sqrt{p/n}\,e^{-d|\log(n/p)|}\le1$, giving the same bound.
At negative frequencies $|\log(np)|\ge\log4$, so the error is at most
\[
 \frac{C_{b,d}p^{1/2-d}}{B^4\ell^2}
       \sum_{n\ge2}\frac{(\log n)^2}{n^{1/2+d}}\ll_{b,d}B^{-4}.
\]
For either sign of $u$ the archimedean density multiplied by $|u|$ is
$O((1+|u|)e^{-|u|/2})$. Using
$e^{-|u|/2}\le p^{-1/2}e^{|u-\ell|/2}$ and the decay of $\eta$ bounds
the normalized integral by $O_{b,d}(\int_\R\sinc^4(Bv)\,dv)=O_{b,d}(1/B)$.
The same computation with $e^{|u|/2}\le\sqrt p\,e^{|u-\ell|/2}$ bounds
the normalized pole terms by $O_{b,d}(p/(B\ell))$. Integrating $|g(u)|$ and
$|ug(u)|$ similarly proves the asserted size and Lipschitz bounds.

For localization, put $F(t)=\int K_B(v)e^{itv}\,dv$. The characteristic
function of the sum $V$ of four independent uniform variables on $[-B,B]$
is $\sinc^4(Bv)$. Fourier inversion gives
\[
 F(t)=\frac{2\pi}{\int k}\,\mathbb E[k(t+V)],\qquad
 h(t)=e^{it\ell}\{F(t-U)+F'(t-U)/(i\ell)\}.
\]
Both $|k(v)|$ and $|k'(v)|$ are $\ll_b\exp(-c_b e^{|v|/b})$. In the two
unit intervals with $H+j<|t-U|\le H+j+1$, $j=0,1,\ldots$, the standard
local zero count is $O(\log(eQ(3+|U|+H+j)))$. Throughout these intervals
$|t-U+V|\ge \Delta+j$. Summation in $j$ gives the stated tail estimate (decreasing
$c_b$ if necessary). To see \eqref{eq:kernel-tail-choice}, first take $\Delta$
at its displayed lower bound. Then $\log(eQ(3+|U|+H))$ is bounded by a
constant times $\log(e^eQ(3+|U|+B+p))$ and increasing $C_{b,d,\epsilon}$
absorbs this constant and $\sqrt p/\ell$. For larger $\Delta$ the omitted
absolute tail can only decrease.
\end{proof}


\section{Averaging the error over the probe primes}
Fix $b>1/\pi$ and fix $d=(1/2+\pi b/2)/2$. The detector targets the prime
$p$ on the logarithmic scale. At $p\asymp x$, a bandwidth $B$ leaves nearby
prime powers within a distance of order $x/B$ on the integer scale. Taking
$B$ proportional to $x/\log x$ lets this distance grow with $x$.
\begin{lemma}[Mean leakage]\label{lem:leak}
Let $1\le D\le\log x$, $B=x/D$. For primes $x<p<2x$ define
\[
 J_p=\sum_{\substack{x/e\le n\le 2ex,\ n\ne p\\ \Lambda(n)>0}}
              (1+|n-p|/D)^{-4}.
\]
There is an absolute $C$ such that for all sufficiently large $x$
\[
 \sum_{x<p<2x} J_p
 \le C\frac{xD}{\log^2 x}+C\sqrt{x}\log(2x)D.
\]
Furthermore the near-prime-power term in Lemma~\ref{lem:kernel-decoder}
is $O(J_p)$ with an absolute constant.
\end{lemma}
\begin{proof}
For the last assertion use the mean value theorem on $[p/e,ep]$ and
$p\asymp x$, so that $|\log(n/p)|\asymp|n-p|/x$ and $\log n/\log p\ll1$.
Put $S(m)=\prod_{r\mid m,\ r>2}(1+1/(r-2))$, with $r$ prime.
For $n$ prime and $h=n-p\ne0$ the two-prime upper-bound sieve gives
\[
 \#\{x<p<2x:\ p+h\text{ prime},\ x/e\le p+h\le2ex\}
 \ll \frac{x}{\log^2 x} S(|h|).
\]
This is uniform in $|h|\le 6x$. For completeness apply the sieve upper
bound \cite[Lemma 3.2]{sieve} to the forbidden residue classes $0,-h$
at primes $r\le\sqrt x$. Both tested primes exceed $\sqrt x$.
For odd $h$ the set is empty. For even $h$, Mertens' product estimate gives
$O(1/\log^2 x)$ from the two forbidden classes, with factor
$(r-1)/(r-2)$ when they coincide at an odd prime. The single factor at $2$
is absorbed in the constant. Expansion over squarefree divisors gives
\[
 \sum_{m\le Y}S(m)\le Y\prod_{r>2\ {\rm prime}}
       \left(1+\frac{1}{r(r-2)}\right)\ll Y.
\]
Summing over dyadic ranges in $|h|/D$ shows that
$\sum_{h\ge1}(1+h/D)^{-4} S(h)\ll D$.
There are $O(\sqrt x\log(2x))$ proper prime powers up to $2ex$, using
at most $O(\log x)$ exponents. For each such $n$ the sum of the weights over
all integer $p$ is $O(D)$ by comparison with an integral. This proves the bound.
\end{proof}

\begin{proof}[Proof of Theorem~\ref{thm:main}]
Take $D=\log x/K$, $B=Kx/\log x$. First choose $K$ large in terms of $b$.
Then choose the fixed $A\ge A_0$ large enough that $D\ge1$ for every $Q\ge3$
and all estimates below hold. Lemma~\ref{lem:leak} and Markov's inequality show
that outside at most $(c_0/2)x/\log x$ primes in $(x,2x)$, the contribution
of the near-prime-power term to the decoder error is at most $1/32$.
In detail, the mean bound divided by $x/\log x$ is
$O(1/K)+O(\log^3(2x)/\sqrt{x})$, and the constants in the decoder are fixed.
Call the primes with $J_p$ below this fixed threshold good. Goodness depends
only on prime powers, $x$ and $b$.

For large $K,A$, the terms $p/B^4$, $1/B$ and the possible pole error
$p/(B\log p)$ in Lemma~\ref{lem:kernel-decoder} contribute at most $1/32$
in total. In particular $p/(B\log p)\le 2/K$. Choose the tail tolerance
as $1/32$. Equation~\eqref{eq:kernel-tail-choice} is satisfied for every
$x<p<2x$ if \eqref{eq:H} holds with $C_b$ large enough.
Indeed $B\asymp_b W(Q)$, and the displayed logarithm in
\eqref{eq:kernel-tail-choice} is at most
$O_b(1+\log\log(3x))+\Lthree(|U|)$. This follows by bounding
$\log(e^e Q(3+|U|+B+p))$ by
$O_b(\log(2Q)+\log(3+|U|)+\log(2x))$ and using
$\log(a+c)\le\log 2+\max(\log a,\log c)$ twice. Consequently,
on each good prime,
\begin{equation}\label{eq:accuracy}
 |D_{p,B,U,H}(\chi)-\chi(p)|<1/8
      \quad\text{simultaneously for every }\chi\in\FQ.
\end{equation}
By Lemma~\ref{lem:witness} at least $(c_0/2)x/\log x$ good primes
separate any prescribed pair by at least one. At any of them the two
spectral measurements differ by more than $3/4$. Identical multisets
would give identical measurements.

For positive heights fix $b=1$, choose $U=KW(Q)$ with a new sufficiently
large absolute constant $K$, and use the window of half-length $U/2$.
The estimate $\Lthree(U)=o(U)$ and increasing $K$ uniformly for $Q\ge3$
give \eqref{eq:H}. This whole window lies in $(0,3KW(Q)/2]$.
\end{proof}

\section{How different must the two observed lists be?}
The proof supplies a quantitative answer even if recorded zeros are perturbed or
omitted. For two finite multisets $S_1,S_2$ in the observation window, pair any
subsets of their copies one-to-one. Let $N_*$ count all unpaired copies on both
sides. Define the cost of this partial matching by
\[
 \mathcal C_x=N_*+
    \sum_{(s,t)\,\mathrm{paired}}\min\{1,(\log x)|s-t|\}.
\]
The cap accounts for the option of treating an appreciably displaced zero as
unmatched at comparable cost.
\begin{corollary}[Separation under partial matching]\label{cor:transport}
With the choices in the proof and $H$ as in \eqref{eq:H}, suppose
$\chi_1\ne\chi_2$. Every partial matching of their lists satisfies
$\mathcal C_x\gg_b\sqrt x\gg_b\log(2Q)$. In particular
\[
 \sum_{t\in[U-H,U+H]} |m_{\chi_1}(t)-m_{\chi_2}(t)|
       \gg_b\log(2Q),
\]
where only distinct ordinates in the union are summed. If the lists have equal
cardinality, the sum of displacements of any bijection is
$\gg_b \log(2Q)/\log\log(3Q)$.
\end{corollary}
\begin{proof}
Use a good witness prime $p$ for the pair. Its two normalized spectral sums differ
by more than $3/4$. Each summand has supremum norm $O_b(\sqrt{x}/(B\log x))$.
Its Lipschitz constant is $O_b(\sqrt{x}/B)$. By the triangle inequality a paired
copy contributes at most the smaller of twice the supremum bound and the
Lipschitz bound times its displacement. An unmatched copy contributes at most
the supremum bound. Thus
$3/4\ll_b\sqrt{x}\,\mathcal C_x/(B\log x)$, and $B\log x=Kx$.
Pair identical copies wherever possible to obtain the first displayed consequence.
For a bijection use $\mathcal C_x\le(\log x)\sum |s-t|$.
\end{proof}

\section{Recovering the character with finite precision}
An individual high-order root of unity need not be distinguishable from every
other character value by an error of $1/8$. Recovery instead compares candidate
characters on a common list of primes. The same list works at every window centre.
\begin{corollary}[A finite alphabet]\label{cor:code}
Fix $b>1/\pi$. For each $Q$ there is a list of $m=O_b(\log(2Q))$ good primes,
allowing repetitions, for which quantizing each spectral measurement to a
fixed finite alphabet identifies $\chi$. An absolute positive fraction of these
symbols can be arbitrarily corrupted, after fixing the universal smooth weight
in Lemma~\ref{lem:witness}. For any fixed alphabet of at least two symbols, the
order $\log Q$ is necessary for identification of the whole family.
\end{corollary}
\begin{proof}
Good witness primes for each pair comprise at least $c_0 x/(2\log x)$ primes.
There are $O(x/\log x)$ primes from which to choose, by the prime number theorem.
Thus a uniformly chosen good prime has probability at least a fixed $\delta>0$
to separate the pair by $1$. Choose $m$ good primes independently. The elementary
binomial lower-tail bound gives probability at most $e^{-\delta m/8}$ of fewer
than $\delta m/2$ witnesses for a given pair. There are at most $Q^2$
characters, since $\sum_{q\le Q}\varphi(q)\le Q^2$, and hence at most $Q^4$
pairs. A union bound proves existence for $m$ a sufficiently large multiple of
$\log(2Q)$. Exhaustive comparison of the finite candidates and prime lists gives
a deterministic choice as well.

Round real and imaginary parts of each measurement to the nearest multiple
of $1/8$. The good measurements have modulus $<9/8$, so one may use the grid
in $[-2,2]^2$ as a fixed alphabet. A rounded symbol is within
$1/8+\sqrt2/16<1/4$ of the correct value $\chi(p)$.
Decode by maximizing the number of coordinates whose symbol lies within
$1/4$ of the candidate's exact value. If $k<\delta m/4$ coordinates are
corrupted, the true candidate matches at least $m-k$ positions. Every competitor
fails in at least $\delta m/2-k>k$ positions. For the converse, a prime modulus
between $Q/2$ and $Q$ supplies $\gg Q$ distinct primitive characters for large
$Q$ by Bertrand's postulate. An alphabet of fixed size $s$ must have $s^m\gg Q$.
\end{proof}
The measurement count refers to linear statistics of all zeros in the chosen
window. The existence proof for the common prime list allows exhaustive search
and provides no polynomial-time complexity claim.

\section{Checks, scope and further questions}
For the two primitive quartic characters modulo $5$, assign $\chi_+(2)=i$ and
$\chi_-(2)=-i$. A numerical sign-change scan gives the following positive
ordinates, in increasing order among the sign changes observed.
\begin{center}
\begin{tabular}{lrrr}
\toprule
Character & First observed & Second observed & Third observed\\
\midrule
$\chi_+$ & 6.183578195451 & 8.457229174423 & 12.67494641701\\
$\chi_-$ & 4.132903705213 & 9.442931129729 & 11.28289644158\\
\bottomrule
\end{tabular}
\end{center}
The character tables and functional equation fix the orientation. The supplied
script checks each residual and its reflected conjugate at 35-digit working
precision. Sampling for sign changes does not certify completeness or
multiplicities. In a separate end-to-end calculation the decoder in
Lemma~\ref{lem:kernel-decoder} used the 35 positive ordinates found for each
character up to height $72$. With $p=2$, $b=1$, $B=8$ and $U=H=36$, it returned
$-0.00017486+0.99997429i$ and $-0.00031971-1.00002114i$, respectively.
Both complex errors were below $0.00033$. The proof of the theorem is independent of
the numerical calculations.

The arithmetic check uses exact integers. For $x=60000$, $B=45000$, among
the 5,244 primes in $(x,2x)$ the script certifies a positive-frequency
near-prime-power error at most $1/8$ for 3,598 probes. For each nontrivial
difference of character exponents modulo $127$, at least 1,818 of those
primes separate the character values by $1$. The three primes $61703$,
$60013$, and $60029$ suffice simultaneously for the 125 nonprincipal
characters in that example. This check isolates the sieve step. It does
not specify a numerical value of the uniform cutoff constant in the theorem.

\paragraph{Further questions.}
Can the order of $W(Q)$ be lowered while retaining a fixed accuracy for each
measurement and a uniform guarantee for every conductor? The proof pays for
primes at the GRH scale $\log^2 Q$ and then gains their mean separation;
a smaller uniform cutoff would require another way of using the distinguishing
information. Can one force an ordinate belonging to exactly one of the two
functions in this same conductor-uniform window? Finally, what explicit,
practical $C_b$ follows from numerical choices of the smooth weight and the sieve
constants? These questions are distinct from the finite uniqueness and the
partial-matching bounds proved here.
\begin{thebibliography}{9}\small
\bibitem{lean} L. Sernau and collaborators, \emph{GRH for Dirichlet
$L$-functions, in Lean 4}, readable edition, commit \texttt{72fead102f66}.
\href{https://github.com/lsernau/grh-lean/blob/72fead102f663d5c7a407964a13dc0cf9d977207/grh/src/Statements.lean}{Statements};
\href{https://github.com/lsernau/grh-lean/blob/72fead102f663d5c7a407964a13dc0cf9d977207/grh/src/Main.lean}{exported theorems};
\href{https://github.com/lsernau/grh-lean/blob/72fead102f663d5c7a407964a13dc0cf9d977207/grh/README.md}{reported checks}.
\bibitem{article} V. Reshetnikov, ProveIt collection, \emph{Beyond the Riemann
hypothesis for Dirichlet $L$-functions}, commit \texttt{dec3e4f21630}.
\href{https://github.com/VladimirReshetnikov/ProveIt/blob/dec3e4f216305071186a1cc4c42fe1cdb38df9be/Analysis/GeneralizedRiemannHypothesis/Article/sections/rec.tex}{Recovery section},
\href{https://github.com/VladimirReshetnikov/ProveIt/blob/dec3e4f216305071186a1cc4c42fe1cdb38df9be/Analysis/GeneralizedRiemannHypothesis/Article/sections/zer.tex}{short-interval section},
\href{https://github.com/VladimirReshetnikov/ProveIt/blob/dec3e4f216305071186a1cc4c42fe1cdb38df9be/Analysis/GeneralizedRiemannHypothesis/Article/README.md}{editorial scope}.
\bibitem{weil} A. Chirre, V. J. Pereira J\'unior and D. de Laat,
\emph{Primes in arithmetic progressions and semidefinite programming},
Math. Comp. 90 (2021), 2235--2246, explicit formula (2.1).
\href{https://doi.org/10.1090/mcom/3638}{doi:10.1090/mcom/3638}.
\bibitem{bmor} M. A. Bennett, G. Martin, K. O'Bryant and A. Rechnitzer,
\emph{Counting zeros of Dirichlet $L$-functions}, Math. Comp. 90 (2021),
1455--1482. \href{https://doi.org/10.1090/mcom/3599}{doi:10.1090/mcom/3599}.
\bibitem{sieve} T. Tao and J. Ter\"av\"ainen, \emph{The Hardy--Littlewood--Chowla
conjecture in the presence of a Siegel zero}, J. Lond. Math. Soc. 106 (2022),
3317--3378, Lemma 3.2. \href{https://doi.org/10.1112/jlms.12663}{doi:10.1112/jlms.12663}.
\bibitem{banks} W. D. Banks, \emph{Comparing zeros of distinct Dirichlet
$L$-functions}, Publ. Math. Debrecen 106 (2025), 327--335.
\href{https://arxiv.org/abs/2302.07073}{arXiv:2302.07073}.
\bibitem{fujii} A. Fujii, \emph{On the zeros of Dirichlet $L$-functions (V)},
Acta Arith. 28 (1976), 395--403.
\href{https://matwbn.icm.edu.pl/ksiazki/aa/aa28/aa2847.pdf}{Publisher PDF}.
\bibitem{loftus} W. Banks and K. Loftus, \emph{Unshared zeros of Dirichlet $L$-functions},
2026 preprint. \href{https://arxiv.org/abs/2607.22930}{arXiv:2607.22930}.
\end{thebibliography}
\end{document}
